\section{Conclusiones}

El desarrollo de esta práctica se llevó a cabo gracias al tutorial de Raspberry Pi de la facultad \cite{FI_UNAM_Raspberry_2026} y a las clases del profesor Moisés \cite{Melendez_Reyes_2026}.

\subsection{Hernandez Acosta Mauricio Gabriel}
A través de esta práctica pude familiarizarme de primera mano con la arquitectura y capacidades de la Raspberry Pi Pico, en especial con el microcontrolador RP2040. Trabajar con el IDE Thonny y MicroPython resultó ser una combinación muy amigable y eficiente, ya que simplifica notablemente la sintaxis sin perder el control directo sobre el hardware. Al desarrollar los algoritmos para los distintos ejercicios —desde el parpadeo básico de un LED hasta el manejo coordinado de múltiples salidas— me quedó claro cómo las instrucciones de código se traducen en señales lógicas reales en los pines GPIO. Esta experiencia me ayudó a comprender mejor la importancia de una correcta gestión de los retardos y cómo estructurar el código para que los periféricos respondan exactamente como se espera en un sistema embebido.

\subsection{Perez Olvera Alexis Abraham}
Cumplir con los objetivos de la práctica me permitió ver el potencial que tiene el RP2040 cuando se combina con la agilidad de MicroPython. Lo más interesante fue observar la evolución gradual de los programas: comenzamos con encendidos básicos y terminamos implementando secuencias más elaboradas, como el contador binario de 8 bits y la temporización cruzada de los dos semáforos. Esto demostró que el diseño algorítmico previo es fundamental antes de pasar a la implementación física, pues un pequeño desfase en los tiempos o un mal manejo de los índices en los pines GPIO altera por completo el comportamiento del circuito. En definitiva, la práctica me brindó bases sólidas para entender la manipulación de terminales digitales y la lógica secuencial en microcontroladores.

\subsection{Sierra Garcia Mariana}
El desarrollo de esta práctica fue fundamental para aterrizar los conceptos teóricos sobre cómo un microcontrolador interactúa con su entorno exterior. Mediante el uso de Thonny y MicroPython, logramos explorar a fondo las terminales GPIO del RP2040 y comprobar cómo una serie de ciclos y estructuras de control bien planteadas permiten gobernar el encendido simultáneo, progresivo y alternado de varios actuadores luminosos. Me pareció especialmente formativo el ejercicio de los semáforos y las secuencias de 8 bits, ya que representan escenarios de control y automatización cercanos a la realidad. Esta práctica reforzó mi confianza para programar sobre hardware embebido y me deja una visión clara de cómo diseñar algoritmos eficientes para el control de puertos de entrada y salida.
